\scanpage{50}
% archive sheet number: 49

\mspage{14}
``\uline{Der 6'}''\inserted{(provisoire)} Si $\mathcal X$ est fibré sur

% slanted
\margin{\deleted{Je \unsure{supprime|suppose} \illegible{1}}}

$\underline\Psi^{\mathrm{opp}}=\underline\Phi={}$\diagram{$\left(\begin{tikzcd}[ampersand replacement=\&,row sep=small]
\& a\arrow[dl]\\
1 \&\\
\& b\arrow[ul]
\end{tikzcd}\right)$}

(donc donné par, $\underline Q$-équiv. près, par un diagramme

% slanted
\margin{\uline{Axiome des fibrations}}

\diagram{\begin{tikzcd}[ampersand replacement=\&,row sep=small]
\& X_a\\
X_1\arrow[ur,no head]\arrow[dr] \&\\
\& X_b
\end{tikzcd}}

dans Cat) alors pour $\xi\in\underline D(\mathcal X)$ \deleted{constant} (voir la \unsure{co|ci|c.}) \deleted{\unsure{H}}\ $\psi_!(\gamma)\in\underline D(\underline Q)$ est non seulement cocartésien (sens évident), mais aussi cartésien.

$\underline\Phi\subset\underline Q^a={}$\diagram{\begin{tikzcd}[ampersand replacement=\&,execute at end picture={\draw (Apoint.north) to[out=178,in=80] ([xshift=-4pt]Onepoint.west) to[out=-95,in=180] (Bpoint.south) to[out=0,in=-92] ([xshift=4pt]Apoint.east) to[out=88,in=0] (Apoint.north) -- cycle;}]
\& |[alias=Apoint]|a\arrow[dl] \&\\
|[alias=Onepoint]|1 \& \& 0\arrow[ul]\arrow[dl]\\
\& |[alias=Bpoint]|b\arrow[ul] \&
\end{tikzcd}}

\uline{NB} Der 6' pour $\underline D$ \deleted{signifie} n'est autre que Der 6 pour $\underline D^\circ$. Je ne vois pas comment l'un impliquerait l'autre, en dehors d'un cas exceptionnel où il y aurait \deleted{\unsure{une dualité}}\inserted{(p.ex.)} un foncteur dualisant dans $A=\underline D(e)$, et une correspondance par dualité entre cohomologie et homologie.

\deleted{Finalement}, il semblerait que la bonne relation, qui englobe Der 6 et Der 6bis, soit la suivante :

% slanted
\margin{\uuline{la condition} \deleted{\unsure{apparue|apparente}}}

% slanted
\margin{``L'\uline{axiome} \uuline{spécifique}''}

\uline{Der\unsure{+|.}6\deleted{\unsure{bis|provisoire}}} Si $\gamma\in\underline D(\Delta^1\times\Delta^1)$, $\gamma$ est cartésien ssi il est cocartésien. \uline{De plus} (\uline{à prouver} que ça en résulte) si diagramme squelettique associé est
